\documentclass[aspectratio=169,professionalfonts]{beamer}

% Compile with XeLaTeX.
% Example:
%   xelatex slides/template/course_beamer_template.tex

\usetheme{Madrid}
\usefonttheme{serif}

\usepackage{fontspec}
\usepackage{xeCJK}

% English fonts
\IfFontExistsTF{TeX Gyre Pagella}{\setmainfont{TeX Gyre Pagella}}{}
\IfFontExistsTF{TeX Gyre Heros}{\setsansfont{TeX Gyre Heros}}{}
\IfFontExistsTF{Menlo}{\setmonofont{Menlo}}{}

% Chinese fonts
% Body text: Kai-style fonts
\IfFontExistsTF{Kaiti SC}{\setCJKmainfont{Kaiti SC}}{%
  \IfFontExistsTF{STKaiti}{\setCJKmainfont{STKaiti}}{%
    \IfFontExistsTF{KaiTi}{\setCJKmainfont{KaiTi}}{%
      \IfFontExistsTF{Songti SC}{\setCJKmainfont{Songti SC}}{}%
    }%
  }%
}

% Title text: Xingkai-style fonts; fallback to Kaiti-style fonts if unavailable
\IfFontExistsTF{Xingkai SC}{\setCJKsansfont{Xingkai SC}}{%
  \IfFontExistsTF{STXingkai}{\setCJKsansfont{STXingkai}}{%
    \IfFontExistsTF{Kaiti SC}{\setCJKsansfont{Kaiti SC}}{%
      \IfFontExistsTF{STKaiti}{\setCJKsansfont{STKaiti}}{}%
    }%
  }%
}

% Avoid xeCJK warnings for monospaced CJK fallback
\IfFontExistsTF{Kaiti SC}{\setCJKmonofont{Kaiti SC}}{%
  \IfFontExistsTF{STKaiti}{\setCJKmonofont{STKaiti}}{%
    \IfFontExistsTF{Songti SC}{\setCJKmonofont{Songti SC}}{}%
  }%
}

% Let title-like elements use the sans family, so Chinese titles follow the Xingkai setting.
% Xingkai fonts often do not provide a true bold face, so keep medium weight here.
\setbeamerfont{title}{family=\sffamily,series=\mdseries}
\setbeamerfont{subtitle}{family=\sffamily,series=\mdseries}
\setbeamerfont{frametitle}{family=\sffamily,series=\mdseries}
\setbeamerfont{section title}{family=\sffamily,series=\mdseries}
\usepackage{amsmath, amssymb, bm}
\usepackage{booktabs}
\usepackage{graphicx}
\usepackage{array}
\usepackage{tikz}
\usepackage{appendixnumberbeamer}

% ==============================
% Metadata: replace these fields
% ==============================
\title[回归分析]{第X周：课件主标题}
\subtitle{Regression Analysis 2026}
\author[课程组]{课程组 / Instructor Name}
\institute[OUC]{中国海洋大学 \quad Regression Analysis 2026}
\date{2026-xx-xx}

% Optional short macros for reuse
\newcommand{\WeekNo}{12}
\newcommand{\TopicName}{Bias-Variance Tradeoff}
\newcommand{\CourseName}{Regression Analysis 2026}

% Math shortcuts
\newcommand{\R}{\mathbb{R}}
\newcommand{\E}{\mathbb{E}}
\newcommand{\Var}{\mathrm{Var}}
\newcommand{\Bias}{\mathrm{Bias}}
\newcommand{\argmin}{\mathop{\mathrm{argmin}}}
\newcommand{\lossfn}{\mathcal{L}}
\newcommand{\regterm}{\Omega}

% Uncomment if needed
% \logo{\includegraphics[height=0.6cm]{figures/logo.pdf}}

\begin{document}

% ------------------------------
% Title page
% ------------------------------
\begin{frame}[plain]
    \titlepage
\end{frame}

% ------------------------------
% Agenda
% ------------------------------
\begin{frame}{本讲导航}
    \tableofcontents
\end{frame}

% ==============================
\section{学习目标}
% ==============================

\begin{frame}{学习目标}
    \begin{itemize}
        \item 理解本讲的核心概念、适用场景与限制条件。
        \item 能够把公式、图像直觉和业务解释连接起来。
        \item 能够根据任务目标选择合适的方法、损失函数或评价指标。
    \end{itemize}
\end{frame}

% ==============================
\section{概念导入}
% ==============================

\begin{frame}{问题背景}
    \begin{block}{典型问题}
        在这里用 2--3 句话说明：为什么需要本讲主题？它解决了什么问题？
    \end{block}

    \vspace{0.4em}
    \begin{itemize}
        \item 背景：给出真实数据分析或机器学习中的触发场景。
        \item 症状：说明传统 OLS / baseline 方法为什么不够。
        \item 目标：本讲希望学生学会什么判断或方法。
    \end{itemize}
\end{frame}

\begin{frame}{一句话主结论}
    \begin{alertblock}{Takeaway}
        在这里放本讲最关键的一句话结论。建议这一页只保留一个中心判断。
    \end{alertblock}

    \vspace{0.5em}
    \begin{exampleblock}{示例写法}
        “当模型复杂度继续上升时，训练误差可能下降，但测试误差不一定继续下降。”
    \end{exampleblock}
\end{frame}

% ==============================
\section{核心理论}
% ==============================

\begin{frame}{核心公式模板}
    \begin{block}{统一目标函数}
        \[
            \hat{\beta} = \argmin_{\beta} \; \lossfn(y, X\beta) + \lambda \, \regterm(\beta)
        \]
    \end{block}

    \begin{itemize}
        \item $\lossfn(\cdot)$：替换为本讲的损失函数，例如平方损失或对数似然。
        \item $\regterm(\cdot)$：替换为本讲的复杂度控制项，例如 $L_1$ 或 $L_2$。
        \item $\lambda$：控制拟合能力与约束强度的超参数。
    \end{itemize}
\end{frame}

\begin{frame}{公式解释模板}
    \begin{columns}[T,onlytextwidth]
        \column{0.48\textwidth}
        \begin{block}{数学层面}
            \begin{itemize}
                \item 定义变量与符号。
                \item 说明每一项的作用。
                \item 指出极端情况：$\lambda = 0$、$\lambda \to \infty$ 等。
            \end{itemize}
        \end{block}

        \column{0.48\textwidth}
        \begin{block}{直观层面}
            \begin{itemize}
                \item 解释“模型在做什么”。
                \item 解释为什么这个公式会影响泛化或可解释性。
                \item 说明和上周知识的连接。
            \end{itemize}
        \end{block}
    \end{columns}
\end{frame}

\begin{frame}{对比表模板}
    \small
    \begin{table}
        \centering
        \begin{tabular}{>{\raggedright\arraybackslash}p{0.18\linewidth} >{\raggedright\arraybackslash}p{0.22\linewidth} >{\raggedright\arraybackslash}p{0.24\linewidth} >{\raggedright\arraybackslash}p{0.24\linewidth}}
            \toprule
            方法 & 核心思想 & 优点 & 风险 / 局限 \\
            \midrule
            方法 A & 在这里填写 & 在这里填写 & 在这里填写 \\
            方法 B & 在这里填写 & 在这里填写 & 在这里填写 \\
            方法 C & 在这里填写 & 在这里填写 & 在这里填写 \\
            \bottomrule
        \end{tabular}
    \end{table}
\end{frame}

\begin{frame}{图示占位模板}
    \begin{columns}[c,onlytextwidth]
        \column{0.58\textwidth}
        \fbox{%
            \parbox[c][0.50\textheight][c]{0.95\linewidth}{
                \centering
                在此插入图像、流程图或实验曲线\\[0.6em]
                推荐内容：bias-variance 曲线、regularization path、PCA 投影示意图等
            }
        }

        \column{0.38\textwidth}
        \begin{block}{讲解提示}
            \begin{itemize}
                \item 图中横轴是什么？
                \item 图中纵轴是什么？
                \item 学生应从图里读出什么结论？
            \end{itemize}
        \end{block}
    \end{columns}
\end{frame}

% ==============================
\section{案例与实践}
% ==============================

\begin{frame}{案例分析模板}
    \begin{block}{案例描述}
        用 2--4 句话描述数据来源、业务背景、目标变量和关键特征。
    \end{block}

    \vspace{0.4em}
    \begin{enumerate}
        \item 问题类型：回归 / 分类 / 计数 / 高维分析。
        \item 主要挑战：异常值、共线性、稀疏性、数据泄露等。
        \item 本讲方法：为什么适合这个场景？
    \end{enumerate}
\end{frame}

\begin{frame}[fragile]{代码模板（可选）}
\small
\begin{verbatim}
from sklearn.pipeline import Pipeline
from sklearn.preprocessing import StandardScaler
from sklearn.linear_model import Ridge
from sklearn.model_selection import GridSearchCV

pipe = Pipeline([
    ("scaler", StandardScaler()),
    ("model", Ridge())
])

param_grid = {"model__alpha": [0.1, 1.0, 10.0]}
search = GridSearchCV(pipe, param_grid=param_grid, cv=5)
search.fit(X_train, y_train)
\end{verbatim}
\end{frame}

\begin{frame}{课堂讨论模板}
    \begin{block}{讨论题}
        在这里放一到两个值得讨论的问题，优先选择“有取舍”的问题，而不是单纯记忆题。
    \end{block}

    \begin{itemize}
        \item 示例 1：什么时候应该优先选择更稳定但更有偏的模型？
        \item 示例 2：在高维问题中，变量筛选和降维各自更适合什么目标？
        \item 示例 3：逻辑回归为什么在工业界仍然具有长期生命力？
    \end{itemize}
\end{frame}

% ==============================
\section{总结与作业}
% ==============================

\begin{frame}{本讲小结}
    \begin{enumerate}
        \item 结论 1：在这里填写本讲的第一条关键结论。
        \item 结论 2：在这里填写本讲的第二条关键结论。
        \item 结论 3：在这里填写本讲与后续课程的连接点。
    \end{enumerate}
\end{frame}

\begin{frame}{作业 / 课后思考}
    \begin{itemize}
        \item 练习题 1：针对本讲核心公式做推导或解释。
        \item 练习题 2：在一个小数据集上完成实验并比较结果。
        \item 练习题 3：用一句业务语言解释模型输出或误差指标。
    \end{itemize}
\end{frame}

\begin{frame}[plain]
    \centering
    \vspace{1.5cm}
    {\usebeamerfont{title}\Huge Q\&A}

    \vspace{1cm}
    {\Large Thanks}
\end{frame}

\end{document}
